\documentclass[11pt]{article}
\usepackage[margin=0.95in]{geometry}
\usepackage{fontspec}
\usepackage{setspace}
\usepackage{hyperref}
\hypersetup{colorlinks=true, allcolors=black}
\setlength{\parindent}{0pt}
\setlength{\parskip}{0.38em}
\setstretch{1.06}
\pagestyle{empty}
\setmainfont{Times New Roman}

\begin{document}

\begin{center}
{\Large Teaching Philosophy}\\[0.25em]
Ke Cheng HSU\\
{\small Department of Economics, The Hong Kong University of Science and Technology}\\
{\small National Chung Cheng University}
\end{center}

My teaching philosophy centers on intuition, active reasoning, and clear explanation. I want students to use economic models not simply to solve equations, but to distinguish an observed pattern from the mechanism that produces it. Mathematics provides the precision needed to state assumptions and trace their implications, yet formulas are meaningful only when students understand the economic question they answer. In every course, my goal is to build the economic intuition first, break complex technical arguments into clear, logical steps, and create an approachable environment where students actively test and refine their understanding.

Three practices guide how I teach. First, I open each topic with a concrete market observation or policy question before introducing formal definitions and equations. Second, I connect intuition to formalism step by step, demonstrating how to set up and solve a model before asking students to practice on an exercise of their own. Third, I treat questions in class and in office hours as diagnostic cues that reveal which step of an argument needs a simpler, more transparent explanation. Teaching is most fulfilling when a student who was stuck begins to see why an economic outcome occurs, and finding that clearer explanation is the part of teaching I enjoy most.

\section*{Teaching Experience at HKUST}

I have taught at the Hong Kong University of Science and Technology (HKUST) in several roles across undergraduate, master's, and doctoral programs. From 2020 to 2025, I served as a teaching assistant across five appointments in undergraduate Macroeconomics, the Economics of Uncertainty and Information, and Ph.D.\ Microeconomic Theory I. In these roles, I led weekly tutorials, held regular office hours, graded assignments and exams, and designed tutorial notes and examination questions. My average evaluation across the five appointments was 4.72 out of 5, including scores of 4.83 in undergraduate information economics and 4.85 and 4.82 in the doctoral microeconomic theory sequence.

Working closely with students across different levels shaped how I explain technical material. Through office hours, I learned to identify the exact point where a student's reasoning becomes unclear and rebuild the argument from there rather than repeating the same derivation. In the doctoral microeconomics course, I often approached the same question in several ways, combining verbal intuition, graphical geometry, and formal algebra until the underlying idea became clear.

In 2026, I served as the instructor for the Mathematics Camp for the MSc in Economics program. I designed the lecture material myself and covered five topics in five days: real analysis, linear algebra, statistics, optimization, and dynamic programming. Because most students entered the program from a non-mathematics background, I built the mathematical foundations step by step and kept the class encouraging. I relied on examples and graphs whenever they could carry the idea, such as using the geometry of projection to explain ordinary least squares before deriving the matrix formula. I received an evaluation of 4.53 out of 5, and students commented that ``the teacher teaches well; let him teach more.'' From 2025 to 2026, I also served as Teaching Assistant Coordinator for the HKUST Business School, helping postgraduate students design teaching materials and evaluating their mock teaching.

\section*{In the Classroom: Intuition Before Formalism}

In my courses, I open each topic with a familiar economic question so students see why a model is needed before they meet the notation. A new restaurant opening on campus illustrates how I introduce information economics. Students may treat a long line outside the restaurant as evidence that the food is good, yet the line is only a signal. As more people observe the line and join, it can keep growing even when the food is bad. This familiar example gives students an intuitive way into herding under imperfect information. A second example, from macroeconomics, asks why a government might issue consumption vouchers instead of cash. Because vouchers must be spent rather than saved, more of the transfer enters current demand and the multiplier is larger. I begin with the intuition of the mechanism and introduce the formula only after students understand the question it answers.

My research in industrial organization and digital economics provides additional applications from everyday markets such as online shopping and retail bookstores. When teaching pricing under uncertainty, for example, I use online purchases and product returns to show how buyer uncertainty shapes a seller's pricing incentive. A return policy affects two margins: whether shoppers buy and whether they return the item after delivery. A stricter return rule cuts wasteful returns, yet lower return rates make a higher price more profitable for the seller and can shut out buyers who would have kept the item. Similarly, when teaching consumer search and product differentiation, we examine how consumers with low search costs seek niche products while high-search-cost consumers buy standardized products, and why a fall in search costs can shift the composition of searchers and raise prices for both product types.

Once the intuition is in place, I break the formal derivation into small, ordered steps and work through a live example before giving students a short in-class exercise. As students work through the problem, I ask them to point out where the reasoning stops being clear. I encourage them to stop me whenever an assumption or step is uncertain, because one student's question often helps several classmates who share the same doubt and shows me where to explain more simply. If an explanation survives their questions in plain language, it usually survives future classes too.

\section*{Course Design and Assessment}

I structure courses around concrete learning outcomes and align lectures, problem sessions, and assessments so that steady practice builds independent reasoning. Rather than treating problem sets as mechanical algebra drills, I design assignments that ask students to set up the optimization or equilibrium problem, solve for the endogenous variables, sign a comparative static, and state which assumption drives the result. In weekly problem sessions, students work through exercises actively, either in pairs at the introductory level or by presenting problem-set solutions in advanced undergraduate and master's courses.

To support students with uneven mathematical preparation while keeping standards high, I pair early review materials with transparent grading. In courses that rely on constrained optimization, the first problem set reviews the Lagrange method before applying it to consumer and producer theory. Students may discuss solution strategies in small groups, while each student writes and submits an independent solution they can reproduce on their own. When grading, I value economic formulation alongside algebraic accuracy: a serious attempt at setting up the model earns credit even when the last algebra step is wrong, whereas a calculation without a clear economic interpretation is incomplete. Examinations use the same structure as the problem sets so that working through the assignments is the direct way to prepare.

\section*{Courses Prepared to Teach and Student Advising}

At National Chung Cheng University, I can teach in both English and Mandarin across the undergraduate and graduate curricula. I am prepared to teach Microeconomic Theory and Industrial Organization from the undergraduate to the Ph.D.\ level, as well as Mathematics for Economists and the Economics of Uncertainty and Information through the master's level. At the undergraduate level, I can also teach Principles of Macroeconomics, Intermediate Macroeconomics, and Econometrics; at the graduate level, I can offer Applied Microeconometrics and Empirical Industrial Organization.

Beyond these theoretical and econometrics courses, I am capable of teaching an introductory course in computational methods as well as a course on conducting research with AI and managing workflows efficiently. I have substantial experience in structural estimation using R, Python, and C++, with Git and GitHub for version control and collaboration. In my research, the structural estimation code was developed in R and C++, while Python supported the data collection and AI-based optical character recognition (OCR) workflows that digitize scanned and historical sources. I treat Git as an essential skill and teach it through a real project so students keep their collaborative work reproducible. I also use AI tools carefully to improve efficiency in coding, debugging, and day-to-day workflow while keeping full ownership of the economic reasoning.

Finally, I look forward to advising undergraduate and graduate student research. In supervising undergraduate projects and master's or doctoral theses, I give tailored feedback at every stage, from framing a feasible economic question, through model building and data work, to the final presentation. Whether in the classroom or in one-on-one advising, my goal is for students to connect evidence to economic mechanisms and build the next step of an argument on their own.

\end{document}
